\documentclass[UTF8,fontset=fandol,10pt,aspectratio=169]{ctexbeamer}

\usepackage{graphicx}
\usepackage{mathtools}
\usepackage{utopia}
\usepackage{booktabs}
\usepackage{xcolor}
\usepackage{natbib}
\usepackage{hyperref}

\usetheme{CambridgeUS}
\usecolortheme{dolphin}

\definecolor{myNewColorA}{RGB}{118,193,188}
\definecolor{myNewColorB}{RGB}{106,172,150}
\definecolor{myNewColorC}{RGB}{94,150,218}
\setbeamercolor*{palette primary}{bg=myNewColorC}
\setbeamercolor*{palette secondary}{bg=myNewColorB, fg=white}
\setbeamercolor*{palette tertiary}{bg=myNewColorA, fg=white}
\setbeamercolor*{titlelike}{fg=myNewColorA}
\setbeamercolor*{title}{bg=myNewColorA, fg=white}
\setbeamercolor*{item}{fg=myNewColorA}
\setbeamercolor*{caption name}{fg=myNewColorA}
\usefonttheme{professionalfonts}
\setbeamertemplate{navigation symbols}{}
\graphicspath{{../../charts/}{../shared/figures/}}

\title[The University of Hong Kong]{数据准备与描述性分析}
\subtitle{JFK 航班延误与取消运营风险项目}
\author[DASE7337 Group Project]{DASE7337 Group Project}
\institute[HKU]{The University of Hong Kong}
\date[April 2026]{DASE7337 Group Presentation \\ April 2026}
\titlegraphic{
\includegraphics[height=1.5cm]{hkulogo.png}}

\setbeamerfont{title}{size=\large}
\setbeamerfont{subtitle}{size=\small}
\setbeamerfont{author}{size=\small}
\setbeamerfont{date}{size=\small}
\setbeamerfont{institute}{size=\small}

\begin{document}

\begin{frame}
  \titlepage
\end{frame}

\begin{frame}{本章在小组报告中的位置}
  % Speaking cue: 30--35 seconds.
  \begin{itemize}
    \item 前一章已经完成行业背景、核心风险因素、SWOT 与 PEST 分析。
    \item 本章负责把 BTS 原始数据整理成可建模的风险数据资产。
    \item 核心问题不是“画了哪些图”，而是“哪些数据口径能支撑后续模型”。
  \end{itemize}
  \vspace{0.4em}
  \begin{block}{建模交接目标}
  将航司月份记录拆分为频率变量、严重度变量和风险原因变量，为后续分布拟合、状态划分与 aggregate risk 模拟提供输入。
  \end{block}
\end{frame}

\begin{frame}{数据范围与清洗结果}
  % Speaking cue: 35--40 seconds.
  \begin{columns}[T,totalwidth=\linewidth]
    \begin{column}{0.48\linewidth}
      \begin{table}
      \small
      \begin{tabular}{lr}
      \toprule
      指标 & 数值 \\
      \midrule
      原始记录 & 570 \\
      清洗后记录 & 569 \\
      覆盖月份 & 72 \\
      覆盖航司 & 12 \\
      总到达航班 & 657,671 \\
      总延误分钟 & 11,125,838 \\
      \bottomrule
      \end{tabular}
      \end{table}
    \end{column}
    \begin{column}{0.48\linewidth}
      \begin{itemize}
        \item 删除 1 条无实际到达航班量的记录。
        \item 检查“延误航班数超过总航班数”的逻辑异常，当前样本为 0。
        \item 统一航司名称，并把 BTS 字段改写为可读变量名。
      \end{itemize}
    \end{column}
  \end{columns}
\end{frame}

\begin{frame}{变量设计：服务 frequency--severity 模型}
  % Speaking cue: 40 seconds.
  \begin{columns}[T,totalwidth=\linewidth]
    \begin{column}{0.33\linewidth}
      \textbf{频率变量}
      \begin{itemize}
        \item 总到达航班量
        \item 15 分钟以上延误
        \item 取消与改降
      \end{itemize}
    \end{column}
    \begin{column}{0.33\linewidth}
      \textbf{严重度变量}
      \begin{itemize}
        \item 总延误分钟
        \item 平均每次延误分钟
        \item 分原因延误分钟
      \end{itemize}
    \end{column}
    \begin{column}{0.33\linewidth}
      \textbf{风险分类变量}
      \begin{itemize}
        \item Airline 与 Late Aircraft
        \item NAS
        \item Weather 与 Security
      \end{itemize}
    \end{column}
  \end{columns}
  \vspace{0.8em}
  \centering
  \large
  数据结构先拆清楚，后续模型才可以分别回答“发生多频繁”和“发生后多严重”。
\end{frame}

\begin{frame}{月度趋势：延误与取消并不完全同步}
  % Speaking cue: 45 seconds.
  \centering
  \includegraphics[width=0.88\linewidth]{chart_1_monthly_trend.png}
  \vspace{-0.3em}
  \begin{itemize}
    \item 最高延误月份：2023 年 7 月，3,657 次延误，延误率 32.69\%。
    \item 最高取消月份：2020 年 3 月，1,900 次取消，体现疫情初期冲击。
  \end{itemize}
\end{frame}

\begin{frame}{风险原因结构：内部链条是主导来源}
  % Speaking cue: 45 seconds.
  \centering
  \includegraphics[width=0.86\linewidth]{chart_2_delay_cause_breakdown.png}
  \vfill
  \begin{minipage}{0.86\linewidth}
    \centering\small
    Airline 与 Late Aircraft 合计贡献约 71.3\% 的延误分钟，说明日常运营风险的主体不是单纯天气，而是航司内部运行链条和系统拥堵。
  \end{minipage}
\end{frame}

\begin{frame}{航司风险画像与 dashboard 交接}
  % Speaking cue: 45--50 seconds.
  \begin{columns}[T,totalwidth=\linewidth]
    \begin{column}{0.60\linewidth}
      \includegraphics[width=\linewidth]{chart_3_airline_risk_profile.png}
    \end{column}
    \begin{column}{0.36\linewidth}
      \begin{itemize}
        \item JetBlue 的总延误分钟最高。
        \item Hawaiian 的延误率最高。
        \item Envoy Air 的平均延误严重度最高。
        \item 离线 dashboard 可用于展示筛选、交互图表和数据表。
      \end{itemize}
      \vspace{0.3em}
      \href{../../dashboard/jfk_operational_risk_dashboard.html}{打开项目 dashboard}
    \end{column}
  \end{columns}
\end{frame}

\begin{frame}{结论：为什么这部分对后续模型重要}
  % Speaking cue: 35--40 seconds.
  \begin{enumerate}
    \item 清洗后的核心数据让频率模型和严重度模型有明确输入。
    \item 原因分类把 BTS 字段转化为课程中的 operational risk blocks。
    \item 月度时间序列为 normal/disrupted 状态划分和 Markov 转移估计打基础。
    \item 图表和 dashboard 让模型结果可以被业务语言解释，而不是停留在公式层面。
  \end{enumerate}
\end{frame}

\end{document}
